\section*{第 \ref{ch:b2:palaeoenvironments} 章——化石与古环境}

\begin{solution}{exo:b2:palaeoenvironments:1}
(a) $112/135 = 0.830$ 和 $112/150 = 0.747$。(b) 海底的死亡贝壳要经过数年
甚至更长时间才堆积起来：其中既有更早生活在此处的物种，也有偶尔出现或被搬运而来的个体，
因此\omterm{def:b2:palaeoenvironments:assemblages}{死亡组合}平均混合了群落的多代成员。(c) 软体类群，
特别是食碎屑者和肉食性动物，如蠕虫：它们没有骨骼，在埋藏之前就已腐烂。
\end{solution}

\begin{solution}{exo:b2:palaeoenvironments:2}
(a) $5 \times 165 = \qty{825}{ka}$，约 \qty{0.8}{Ma}，处于 0.7 至 1 百万年之间。
(b) $183.73 - 182.77 = \qty{0.96}{Ma}$，$0.96/0.165 = 5.8$：约六个层位。(c) 谱系演化
在远短于一百万年的时间内完成；只有层位能区分相继出现的类型，
而生物带则将它们归并在一起。
\end{solution}

\begin{solution}{exo:b2:palaeoenvironments:3}
(a) Nereites \omterm{def:b2:palaeoenvironments:ichnofacies}{遗迹相}：陆架坡折以外浊积体系的深海海底，
食物稀少。(b) Skolithos \omterm{def:b2:palaeoenvironments:ichnofacies}{遗迹相}：前滨或临滨。(c) Cruziana \omterm{def:b2:palaeoenvironments:ichnofacies}{遗迹相}：
位于近海，介于正常浪基面与风暴浪基面之间。(d) Zoophycos \omterm{def:b2:palaeoenvironments:ichnofacies}{遗迹相}：在
中生代，位于陆架下部或斜坡，风暴浪基面以下。
\end{solution}

\begin{solution}{exo:b2:palaeoenvironments:4}
(a) A 层含狭盐性类群，指示正常海水。B 层物种少而
个体多，含有植物碎屑，且没有任何狭盐性类群，提示为咸淡水环境，例如
河口湾。(b) 缺失的类群有若干种，在同时代的海相地层中常见且保存良好：
它们共同缺失是一种规律，而单个化石的缺失则可能出于偶然或
保存不佳。(c) 沉积特征（潮汐或河道、泥质披覆、根系）、\omterm{def:b2:palaeoenvironments:fossil-kinds}{遗迹化石}（
物种较少的 Cruziana \omterm{def:b2:palaeoenvironments:ichnofacies}{遗迹相}），以及 B 层向侧向过渡为
海相或河流沉积。
\end{solution}

\begin{solution}{exo:b2:palaeoenvironments:5}
(a) $P = 1/2$ 即 $kh/s = \ln 2$，于是 $s = 0.05 \times 0.1/0.693 = \qty{0.0072}{m.\text{年}^{-1}}$，
即每年 \qty{7.2}{mm}。(b) 将 $s$ 减半使指数加倍：$P = (1/2)^2 = 0.25$，为原来的一半。
(c) $k$ 减半：$P = e^{-2.5} = 0.082$，而非 $e^{-5} = 0.0067$，
可能性高出 $e^{2.5} = 12$ 倍。(d) 厚实、坚固的贝壳被过度代表，且埋藏越慢，这种偏差越明显。
\end{solution}

\begin{solution}{exo:b2:palaeoenvironments:6}
(a) 上限为 $183.73 + 0.35 = \qty{184.08}{Ma}$：图表中的 \qty{184.2}{Ma}
超出此上限，相差 \qty{0.12}{Ma}。(b) $183.73 - 182.77 = \qty{0.96}{Ma}$。(c) $6/0.96 =
\qty{6.25}{m}$ 每百万年。(d) 生物带由化石界定，化石的出现与消失
在广阔区域内近乎同时发生；该时段内形成的岩层厚度
则取决于当地的沉积速率。
\end{solution}

\begin{solution}{exo:b2:palaeoenvironments:7}
(a) $0.49\,e^{-1046/3704} = 0.369$ 和 $0.49\,e^{-1166/3704} = 0.358$；平均值为 $0.364$。(b) $120
\times (1 - 0.364) = \qty{76.3}{m}$。(c) $76.3/(1 - 0.49) = \qty{149.6}{m}$。(d) 该因子
假定整个恢复层段都具有地表的孔隙度 $\varphi_0$；但地表厚度为
\qty{150}{m} 的层段，其底部比顶部的孔隙度更低，因此同样的颗粒
填充的厚度小于 \qty{149.6}{m}。
\end{solution}

\begin{solution}{exo:b2:palaeoenvironments:8}
(a) 单位体积中含有 $\varphi$ 的水和 $1 - \varphi$ 的颗粒：$\rho = \varphi
\rho_{\text{水}} + (1 - \varphi)\rho_{\text{颗粒}}$，由此得到该式；$(2710 -
2573)/(2710 - 1000) = 0.080$。(b) $0.51\,e^{-956/4545} = 0.413$。(c) 置换沉淀：方解石
在孔隙中沉淀，碳酸盐岩在压实之外还经历这一过程，其去除的孔隙度
远多于单纯埋藏所造成的。
\end{solution}

\begin{solution}{exo:b2:palaeoenvironments:9}
(a) 深海浊积岩（Nereites \omterm{def:b2:palaeoenvironments:ichnofacies}{遗迹相}）；侏罗纪时位于陆架下部或斜坡的海底
（Zoophycos \omterm{def:b2:palaeoenvironments:ichnofacies}{遗迹相}）；介于浪基面之间的近海海底（Cruziana \omterm{def:b2:palaeoenvironments:ichnofacies}{遗迹相}）；临滨（Skolithos \omterm{def:b2:palaeoenvironments:ichnofacies}{遗迹相}）；
陆上沼泽。(b) 水体向上变浅：海岸向海方向推进。(c) 海平面下降，
或沉积物供应速度超过沉降速度，如同生长中的三角洲前方；若为海平面下降，
则同期其他盆地中也应可见，而三角洲则是局地的，其沉积物来自一侧。
(d) 在整合层序中，相邻的沉积环境依次相接；
若跨越中间的环境，则指示时间上的间断或断层。
\end{solution}

\begin{solution}{exo:b2:palaeoenvironments:10}
(a) 掩埋所需时间为 $h/s$ = 100、10 和 \qty{1}{\text{年}}。软体：$e^{-100} \approx
\num{4e-44}$，$e^{-10} = \num{4.5e-5}$，$e^{-1} = 0.37$。贝壳（$k = \qty{0.01}{\text{年}^{-1}}$）：
$e^{-1} = 0.37$，$e^{-0.1} = 0.90$，$e^{-0.01} = 0.99$。(b) $s = kh/\ln 100 = 0.1/4.605 =
\qty{0.022}{m.\text{年}^{-1}}$，即每年 \qty{22}{mm}。(c) 几乎不会，历时数月；但风暴层、
浊积岩或火山灰沉降可在数小时内掩埋海底。(d) 若沉积速率恒定，
则掩埋所需时间远短于按平均速率算得的 $h/s$，而 $P = e^{-kT}$，其中 $T$
为实际掩埋时间，当 $T$ 远小于 $1/k$ 时趋于 1。
\end{solution}

\begin{solution}{exo:b2:palaeoenvironments:11}
(a) 薄层 $\mathrm dz$ 位于深度 $z$ 处，含有颗粒 $(1 - \varphi(z))\,\mathrm dz$；
在地表，它们占据 $\mathrm dz_0$，且 $(1 - \varphi_0)\,\mathrm dz_0 = (1 -
\varphi(z))\,\mathrm dz$。(b) $\varphi = 0.63\,e^{-816/1960} = 0.415$；$(1 - 0.415)/0.37 =
1.581$，即 \qty{158.1}{m}。(c) 多出 \qty{9.5}{m}：如\cref{exo:b2:palaeoenvironments:7}所示，
恢复的层段并不像其地表那样均匀地致密。(d) 颗粒通量
$h(1-\varphi)/\Delta t$ 仅用到现今的厚度和孔隙度，二者均为实测值。
\end{solution}

\begin{solution}{exo:b2:palaeoenvironments:12}
(a) $9.5/0.165 = 57.6$ 个层位。(b) $180/57.6 = \qty{3.1}{m}$ 为今日的石灰岩，$338/57.6 =
\qty{5.9}{m}$ 为新沉积物。(c) 至多 \qty{1.5}{m}：$180/1.5 = 120$ 个间隔，121
个样品。(d) 沉积间断、菊石缺失的岩层，以及被压缩为
数厘米厚的层位，即沉积物堆积缓慢之处。
\end{solution}

\begin{solution}{pb:b2:palaeoenvironments:1}
\textbf{1.} $866 \sin \ang{12} = 866 \times 0.2079 = \qty{180.0}{m}$。
\textbf{2.} $481 \times 0.2079 = \qty{100.0}{m}$。
\textbf{3.} 据叠覆原理，砂岩最老：两翼倾角仅为 \ang{12}，地层未倒转，
因此最底层者最老。
\textbf{4.} $866 - 766 = \qty{100}{m}$ 和 $1046 - 866 = \qty{180}{m}$。
\textbf{5.} 该处至少为 \qty{766}{m}；其原始顶部已被侵蚀，因此
全厚未知。
\textbf{6.} 地表为 0.70；在 \qty{766}{m} 处为 $0.70\,e^{-766/1408} = 0.406$。
\textbf{7.} 它们是狭盐性类群：水体盐度与开阔海相同。
\textbf{8.} Cruziana \omterm{def:b2:palaeoenvironments:ichnofacies}{遗迹相}：近海，介于正常浪基面与
风暴浪基面之间。
\textbf{9.} 风暴：其波浪到达海底，打碎贝壳，并将其铺成
具侵蚀底面的层。
\textbf{10.} 底层水含氧量低：穿掘生物稀少，纹层得以保存，并形成了黄铁矿。
\textbf{11.} Zoophycos \omterm{def:b2:palaeoenvironments:ichnofacies}{遗迹相}：中生代时位于陆架下部或斜坡，
风暴浪基面以下；有时氧气足以供穿掘生物生存。
\textbf{12.} 否：菊石在上层水体中游动；死后其贝壳下沉或漂移至海底，
形成\omterm{def:b2:palaeoenvironments:assemblages}{死亡组合}。
\textbf{13.} $184.2 - 174.7 = \qty{9.5}{Ma}$。
\textbf{14.} $180/9.5 = \qty{18.9}{m}$ 每百万年。
\textbf{15.} $(2710 - 2573)/(2710 - 1000) = 0.080$。
\textbf{16.} $180 \times (1 - 0.080)/(1 - 0.51) = \qty{338}{m}$。
\textbf{17.} $338/9.5 = \qty{35.6}{m}$ 每百万年。
\textbf{18.} $180 \times 0.92/9.5 = \qty{17.4}{m}$ 每百万年。
\textbf{19.} 页岩段给出 11.3 和 \qty{4.4}{m} 每百万年：颗粒通量下降
约四倍。随着海底加深，陆架与浅滩来源的碳酸盐不再输入，
只有缓慢的泥质沉积物到达此处。
\textbf{20.} 海水向陆地推进：先是河流，随后是浅滩、近海海底、风暴
浪基面以下的海底，最后是白垩所代表的开阔海。该柱状剖面记录了一次\omterm{def:b2:phanerozoic-earth:transgression}{海侵}。
\textbf{21.} 石灰岩之上为 $57 - 41 = \qty{16}{m}$，页岩之上为 $86 - 57 = \qty{29}{m}$，
共计 \qty{45}{m}。
\textbf{22.} $(60 - 5) + 338 + 148.6 = \qty{542}{m}$，历时 $184.2 - 161.5 = \qty{22.7}{Ma}$：
\qty{23.9}{m} 每百万年。
\textbf{23.} $45/542 = 0.083$：约 \qty{8}{\percent}，若旋回图夸大了
振幅，则更少。其余部分来自海底的沉降。
\textbf{24.} 瓦尔登石灰岩的灰泥在沉积时以约
\textbf{\qty[text-series-to-math=true]{36}{m}} 每百万年的速率堆积。
\end{solution}
